% Readout 2: fixed-slot head initialised from lm_head rows.
% Same input/gather as readout 1; head is an independent Linear(d,K).
% W_slot stored [K,d], shown transposed [d,K]; W_slot^(0)=W_lm[ell,:].
% pplx variant: full-attention layers use an all-ones (non-causal) mask.
% MetisCraft readout tensor figures: shared notation, palette and macros.
% Cell counts are illustrative; each symbolic axis keeps one physical edge
% length inside a figure (u per cell): T=12, d=5, V=12, K=L=3, B=2, d_h=4.
\documentclass[tikz,border=3mm]{standalone}
\usepackage[UTF8,fontset=fandol]{ctex}
\usepackage{amsmath,amssymb,bm}
\usetikzlibrary{arrows.meta,calc,backgrounds,decorations.pathreplacing}
\definecolor{Ink}{HTML}{273648}
\definecolor{Muted}{HTML}{607086}
% token roles
\definecolor{Ctx}{HTML}{85898F}    % state / context
\definecolor{Qst}{HTML}{4F8FA5}    % question
\definecolor{Opt}{HTML}{5A9A6E}    % options / candidates (three lightness levels)
\definecolor{Gen}{HTML}{B89A5E}    % generated rationale tokens
% tensor roles
\definecolor{Hid}{HTML}{8A74B5}    % hidden activations
\definecolor{Wgt}{HTML}{E0904F}    % trainable / pretrained weights
\definecolor{Lgt}{HTML}{C95B5B}    % logits, probabilities, loss
\definecolor{Idx}{HTML}{7E8B9B}    % integer indices / labels
\definecolor{Msk}{HTML}{3F5A78}    % binary attention support
\colorlet{TokS}{Ctx!55}\colorlet{TokQ}{Qst!65}
\colorlet{Tok1}{Opt!30}\colorlet{Tok2}{Opt!58}\colorlet{Tok3}{Opt!88}
\colorlet{TokD}{Ink!80}\colorlet{TokG}{Gen!65}\colorlet{TokM}{Opt!100!black!60}
\colorlet{TokC}{Ink!80}\colorlet{TokP}{white}
\colorlet{TokL}{Opt!58}\colorlet{TokT}{Qst!65}
\newcommand{\U}{0.3}
\newcommand{\Off}{0.14}
\newcounter{tseed}
\tikzset{
  flow/.style={-{Latex[length=2mm,width=1.4mm]},draw=Ink!85,line width=.7pt},
  dflow/.style={-{Latex[length=2mm,width=1.4mm]},draw=Ink!60,line width=.6pt,dashed},
  op/.style={draw=Ink!30,fill=white,rounded corners=2pt,line width=.5pt,
             align=center,font=\small,inner sep=4pt},
  sym/.style={font=\small,inner sep=1pt},
  shp/.style={font=\scriptsize,text=Muted,inner sep=1pt},
  note/.style={font=\scriptsize,text=Muted,align=center,inner sep=1pt},
  stage/.style={font=\small\bfseries,text=Muted},
  bigop/.style={font=\Large,text=Ink},
  brc/.style={decorate,decoration={brace,amplitude=3pt},draw=Ink!55,line width=.5pt},
}
% \Tens{id}{cx}{cy}{rows}{cols}{color}{depth}{symbol}{shape}
\newcommand{\Tens}[9]{%
  \stepcounter{tseed}%
  \pgfmathsetmacro{\tw}{#5*\U}\pgfmathsetmacro{\th}{#4*\U}%
  \pgfmathsetmacro{\tx}{#2-\tw/2}\pgfmathsetmacro{\ty}{#3-\th/2}%
  \ifnum#7>0
    \foreach \k in {#7,...,1}{%
      \path[draw=#6!45,fill=white,rounded corners=.8pt,line width=.4pt]
        ({\tx+\k*\Off},{\ty+\k*\Off}) rectangle ++(\tw,\th);}%
  \fi
  \foreach \r in {1,...,#4}{\foreach \c in {1,...,#5}{%
    \pgfmathtruncatemacro{\lv}{mod(int(abs(sin(\r*127.1+\c*311.7+\arabic{tseed}*74.7))*1000),7)}%
    \pgfmathtruncatemacro{\tone}{\lv<3?30:(\lv<6?55:80)}%
    \path[fill=#6!\tone,rounded corners=.6pt]
      ({\tx+(\c-1)*\U+.022},{\ty+\th-\r*\U+.022}) rectangle ++({\U-.044},{\U-.044});}}%
  \coordinate (#1-w) at (\tx,#3);
  \coordinate (#1-e) at ({\tx+\tw+#7*\Off},#3);
  \coordinate (#1-n) at (#2,{\ty+\th+#7*\Off});
  \coordinate (#1-s) at (#2,\ty);
  \coordinate (#1-nw) at (\tx,{\ty+\th});
  \coordinate (#1-sw) at (\tx,\ty);
  \node[sym] (#1-sym) at (#2,{\ty-.27}) {$#8$};
  \node[shp] (#1-shp) at (#2,{\ty-.6}) {$#9$};}
% outline one row (r, 1-based) of a tensor whose front face top-left is (x,y)
% \RowMark{x_left}{y_top}{row}{cols}{color}
\newcommand{\RowMark}[5]{%
  \path[draw=#5,line width=1.1pt,rounded corners=1pt]
    ({#1-.03},{#2-#3*\U-.03}) rectangle ++({#4*\U+.06},{\U+.06});}
\newcommand{\ColMark}[5]{% x_left y_top col rows color
  \path[draw=#5,line width=1.1pt,rounded corners=1pt]
    ({#1+(#3-1)*\U-.03},{#2-#4*\U-.03}) rectangle ++({\U+.06},{#4*\U+.06});}
% \TokRow{x_left}{y_center}{roles}: one tile per token, roles in {S,Q,1,2,3,D,G,M,C,P,L,T}
\newcommand{\TokRow}[3]{%
  \foreach \t [count=\i] in {#3}{%
    \path[fill=Tok\t,draw=Ink!25,line width=.3pt,rounded corners=.6pt]
      ({#1+(\i-1)*\U+.022},{#2-\U/2+.022}) rectangle ++({\U-.044},{\U-.044});}}
\newcommand{\TokCol}[3]{% x_center y_top roles (top to bottom)
  \foreach \t [count=\i] in {#3}{%
    \path[fill=Tok\t,draw=Ink!25,line width=.3pt,rounded corners=.6pt]
      ({#1-\U/2+.022},{#2-\i*\U+.022}) rectangle ++({\U-.044},{\U-.044});}}
% \Mask{x_left}{y_top}{n}{groups}{mode}
%  mode 0: causal & (g_j=0 or g_j=g_i)      shared-prefix tree
%  mode 1: all ones                          bidirectional
%  mode 2: causal & (g_j=0 or g_j=g_i or g_i=9)  pointer (9 = decision token)
%  mode 3: causal                            plain causal
%  mode 4: g_j=g_i                           independent bidirectional blocks
%  mode 5: causal with generated tokens (g=5) drawn in Gen colour
\newcommand{\Mask}[5]{%
  \foreach \i in {1,...,#3}{\foreach \j in {1,...,#3}{%
    \pgfmathtruncatemacro{\gi}{{#4}[\i-1]}%
    \pgfmathtruncatemacro{\gj}{{#4}[\j-1]}%
    \ifcase#5
      \pgfmathtruncatemacro{\ok}{and(\j<=\i,or(\gj==0,\gj==\gi))}\or
      \pgfmathtruncatemacro{\ok}{1}\or
      \pgfmathtruncatemacro{\ok}{and(\j<=\i,or(or(\gj==0,\gj==\gi),\gi==9))}\or
      \pgfmathtruncatemacro{\ok}{\j<=\i}\or
      \pgfmathtruncatemacro{\ok}{\gj==\gi}\or
      \pgfmathtruncatemacro{\ok}{\j<=\i}\fi
    \ifnum\ok=1
      \path[fill=Msk!70,rounded corners=.5pt]
        ({#1+(\j-1)*\U+.022},{#2-\i*\U+.022}) rectangle ++({\U-.044},{\U-.044});
    \fi}}%
  \path[draw=Ink!40,line width=.4pt] (#1,#2) rectangle ++({#3*\U},{-#3*\U});}
% integer slot strip: \Slots{x_left}{y_center}{entries}
\newcommand{\Slots}[3]{%
  \foreach \v [count=\i] in {#3}{%
    \node[draw=Idx!70,fill=Idx!12,rounded corners=1pt,minimum width=\U cm-.04cm,
      minimum height=\U cm-.04cm,inner sep=0pt,font=\tiny,text=Ink]
      at ({#1+(\i-.5)*\U},#2) {$\v$};}}
% bottom meaning box: \Meaning{x0}{y_top}{width}{axes}{objects}{mechanism}{id line}
\newcommand{\Meaning}[7]{%
  \node[anchor=north west,fill=Ink!4,rounded corners=3pt,inner sep=9pt,
    text width=#3cm-18pt,align=left,font=\small] (meaning) at (#1,#2) {%
    \begin{tabular}{@{}p{1.1cm}@{\hspace{6pt}}p{\dimexpr #3cm-2.0cm\relax}@{}}
    \textbf{轴} & \raggedright #4\tabularnewline[5pt]
    \textbf{对象} & \raggedright #5\tabularnewline[5pt]
    \textbf{机制} & \raggedright #6\tabularnewline
    \end{tabular}};
  \node[font=\scriptsize,text=Ink!45,anchor=north] at ($(meaning.south)+(0,-.18)$) {#7};}

\begin{document}
\begin{tikzpicture}[x=1cm,y=1cm,text=Ink]
\node[font=\large] at (10.6,4.05)
  {$W_{\mathrm{slot}}^{(0)}=W_{\mathrm{lm}}[\ell,:],\quad
    z_b=W_{\mathrm{slot}}\,h_b,\quad
    p_b=\operatorname{softmax}_{k}(z_b),\quad
    \mathcal L=-\tfrac1B\textstyle\sum_b\log p_{b,y_b}$};
% ---------- stage A ----------
\node[stage] at (10.3,3.35) {输入与读出 1 相同；mask 默认因果，pplx 把全注意力层改为非因果};
\node[sym,anchor=east] at (1.0,2.3) {$X$};
\node[shp,anchor=east] at (1.0,2.0) {$[B,T]$};
\TokRow{1.2}{2.15}{S,S,S,Q,Q,1,1,2,2,3,3,D}
\draw[brc] (1.22,2.36) -- (2.08,2.36) node[midway,above=3pt,note] {state};
\draw[brc] (2.12,2.36) -- (2.68,2.36) node[midway,above=3pt,note] {问题};
\draw[brc] (2.72,2.36) -- (4.48,2.36) node[midway,above=3pt,note] {选项};
\draw[brc] (4.52,2.36) -- (4.78,2.36) node[midway,above=3pt,note] {$t^\star$};
\Mask{1.2}{1.8}{12}{0,0,0,0,0,0,0,0,0,0,0,0}{3}
\node[sym] at (3.0,-2.07) {$M$};
\node[shp] at (3.0,-2.4) {$[T,T]$ 因果（默认）};
\node[font=\small,text=Muted] at (5.3,0) {或};
\TokRow{5.8}{2.15}{S,S,S,Q,Q,1,1,2,2,3,3,D}
\Mask{5.8}{1.8}{12}{0,0,0,0,0,0,0,0,0,0,0,0}{1}
\node[sym] at (7.6,-2.07) {$M^{\mathrm{bi}}$};
\node[shp] at (7.6,-2.4) {$[T,T]$ 全 1（pplx）};
\node[op,minimum width=2.4cm,minimum height=1.3cm] (dec) at (11.35,0)
  {\textbf{Decoder}\\[2pt]$M$ 或 $M^{\mathrm{bi}}$};
\draw[flow] (9.55,0) -- (dec.west);
\Tens{H}{13.95}{0}{12}{5}{Hid}{1}{H}{[B,T,d]}
\RowMark{13.2}{1.8}{12}{5}{Lgt}
\draw[flow] (dec.east) -- (H-w);
\node[op,minimum width=2.1cm,minimum height=1.1cm] (g1) at (16.75,0)
  {\textbf{Gather}\\[2pt]$H_{b,t^\star_b,:}$};
\draw[flow] (H-e) -- (g1.west);
\Slots{16.45}{-1.25}{11,11}
\node[sym] at (16.75,-1.62) {$t^\star$};
\node[shp] at (16.75,-1.95) {$[B]$};
\draw[dflow] (16.75,-1.08) -- (g1.south);
\Tens{h}{19.4}{0}{2}{5}{Hid}{0}{h}{[B,d]}
\draw[flow] (g1.east) -- (h-w);
% ---------- stage B ----------
\node[stage] at (13.0,-3.3) {独立可训练头 $\mathrm{Linear}(d,K)$：每个槽位一行，沿 $K$ 归一化};
\Tens{h2}{8.0}{-5.0}{2}{5}{Hid}{0}{h}{[B,d]}
\node[bigop] at (9.2,-5.0) {$\times$};
\Tens{Ws}{10.15}{-5.0}{5}{3}{Wgt}{0}{W_{\mathrm{slot}}^{\mathsf T}}{[d,K]}
\Tens{z}{12.3}{-5.0}{2}{3}{Lgt}{0}{z}{[B,K]}
\draw[flow] (Ws-e) -- (z-w);
\node[op,minimum width=1.6cm,minimum height=1.1cm] (sm) at (14.35,-5.0)
  {softmax\\[2pt]沿 $K$};
\draw[flow] (z-e) -- (sm.west);
\Tens{p}{16.4}{-5.0}{2}{3}{Lgt}{0}{p}{[B,K]}
\draw[flow] (sm.east) -- (p-w);
\node[op,minimum width=2.4cm,minimum height=1.1cm] (loss) at (19.05,-5.0)
  {\textbf{交叉熵}\\[2pt]$-\log p_{b,y_b}$};
\draw[flow] (p-e) -- (loss.west);
\Slots{18.75}{-6.15}{0,2}
\node[sym] at (19.05,-6.52) {$y$};
\node[shp] at (19.05,-6.85) {$[B]$};
\draw[dflow] (19.05,-5.98) -- (loss.south);
% ---------- stage C: initialisation ----------
\node[stage] at (6.0,-7.25) {训练前一次：用标签 token 的 lm\_head 行初始化};
\Tens{Wl}{3.2}{-9.0}{5}{12}{Wgt}{0}{W_{\mathrm{lm}}^{\mathsf T}}{[d,V]}
\foreach \c in {4,8,11}{\ColMark{1.4}{-8.25}{\c}{5}{Ink!80}}
\node[op,minimum width=1.9cm,minimum height=1.1cm] (cp) at (7.35,-9.0)
  {\textbf{复制列}\\[2pt]$W_{\mathrm{lm}}^{\mathsf T}[:,\ell]$};
\draw[flow] (Wl-e) -- (cp.west);
\Slots{6.9}{-10.15}{3,7,10}
\node[sym,anchor=west] at (7.85,-10.15) {$\ell$\;{\scriptsize\color{Muted}$[K]$}};
\draw[dflow] (7.35,-9.98) -- (cp.south);
\Tens{W0}{10.15}{-9.0}{5}{3}{Wgt}{0}{W_{\mathrm{slot}}^{(0)\mathsf T}}{[d,K]}
\draw[flow] (cp.east) -- (W0-w);
\draw[dflow] (W0-n) -- (10.15,-6.58);
\node[note,anchor=west] at (10.3,-7.55) {作初值};
% ---------- meaning ----------
\Meaning{0.0}{-11.5}{21.6}
 {$B$ 问题数；$T$ 序列长度；$d$ hidden；$V$ 词表；$K$ 槽位数 = 可表达的选项上限（pplx 与 d3 为 255，AutoTrust 264）。}
 {$W_{\mathrm{slot}}$ 独立可训练权重，存储 $[K,d]$，图中转置；$W_{\mathrm{slot}}^{(0)}$ 其初值；$\ell$ 槽位对应 token 的词表 id；$M$ 因果、$M^{\mathrm{bi}}$ 全 1 的可见性；$z$ logits，$p$ 槽位上的概率。}
 {与读出 1 共用输入和 gather，只把“读 $V$ 列再挑 $L$ 列”换成可训练的 $K$ 行；导出时可把 $W_{\mathrm{slot}}$ 写回 lm\_head 的 $\ell$ 行供普通引擎服务。改用 $M^{\mathrm{bi}}$ 后，普通 CausalLM 推理不再等价。}
 {读出 2：lm\_head 行初始化的固定槽位头@Shichao}
\end{tikzpicture}
\end{document}
